\documentclass[conference]{IEEEtran}
\IEEEoverridecommandlockouts
\usepackage{cite}
\usepackage{amsmath,amssymb,amsfonts}
\usepackage{graphicx}
\usepackage{textcomp}
\usepackage{xcolor}
\usepackage{array}

% Tables and figures are read from tables/ and figures/ next to this file.
% code/main.m, code/make_paper_outputs.m and code/export_to_paper.m create them.
\newcommand{\inputtable}[1]{\IfFileExists{tables/#1.tex}{\input{tables/#1.tex}}{\fbox{\footnotesize missing: tables/\detokenize{#1}.tex}}}
\newcommand{\inputfig}[2]{\IfFileExists{figures/#1.pdf}{\includegraphics[width=#2]{figures/#1.pdf}}{\fbox{\footnotesize missing: figures/\detokenize{#1}.pdf}}}

\begin{document}

\title{Credit Supply, Uncertainty and U.S. Private R\&D Investment: Evidence from a Sign-Restricted VAR}

\author{\IEEEauthorblockN{Albert Moln\'ar}
\IEEEauthorblockA{\textit{Keleti K\'aroly Faculty of Business and Management} \\
\textit{\'Obuda University}\\
Budapest, Hungary \\
molnar.albert@kgk.uni-obuda.hu}
}

\maketitle

\begin{abstract}
This paper studies how credit supply shocks and economic policy uncertainty
shocks affect private research and development (R\&D) investment in the
United States, the largest R\&D economy in the world and the centre of the
current investment cycle in artificial intelligence. We estimate a Bayesian structural vector
autoregression on quarterly data for 1986--2019 and identify credit supply,
uncertainty and aggregate supply shocks with sign and relative-magnitude
restrictions on output, prices, the one-year Treasury yield, the excess bond
premium and the economic policy uncertainty index. The response of R\&D
investment is left unrestricted. A credit supply shock of one standard
deviation lowers real R\&D investment by about 1\% after two years, and the
posterior probability of a decline is close to 0.9. The response to an
uncertainty shock is smaller and its sign is not robust across
specifications. Local projections do not detect a lasting effect on the level
of R\&D investment five to seven years after either shock, although the
shortfall in R\&D investment implies a persistent, if imprecisely estimated,
loss of R\&D capital.
\end{abstract}

\begin{IEEEkeywords}
R\&D investment, credit supply, excess bond premium, policy uncertainty,
structural VAR, sign restrictions, local projections
\end{IEEEkeywords}

% ===========================================================================
\section{Introduction}
% ===========================================================================

Whether business R\&D investment falls in downturns has been debated since
Schumpeter. In the opportunity-cost view of Aghion and Howitt
\cite{aghionhowitt1992}, recessions lower the return to current production
relative to innovation, so firms should shift resources towards R\&D when
activity is weak. Most of the empirical evidence points the other way:
aggregate and industry-level R\&D tends to be procyclical
\cite{rafferty2003,barlevy2007,ouyang2011,fabrizio2014}. Two explanations are
usually offered. R\&D is financed largely from internal funds and equity
because it offers little collateral, so a contraction of credit supply forces
constrained firms to cut it \cite{himmelberg1994,brown2009,aghion2012}.
R\&D projects are also partly irreversible, so higher uncertainty raises the
value of postponing them \cite{bernanke1983,bloom2009}. Because R\&D
accumulates into a stock of knowledge, temporary cuts can lower the level of
productivity for a long time \cite{anzoategui2019,cerra2023}.

The United States is the natural place to study these mechanisms. It
accounted for about 28\% of world R\&D expenditure in 2019, and the business
sector performed about three quarters of U.S. R\&D \cite{nsb2022}. U.S.-based
firms accounted for 42\% of the R\&D spending of the world's 2{,}000 largest
corporate R\&D investors in 2023, and for about 70\% in ICT software
\cite{nindl2024}. The current investment cycle in artificial intelligence is
concentrated in the same economy: private investment in AI reached \$109
billion in the United States in 2024, nearly twelve times the Chinese figure
\cite{maslej2025}. In the national accounts data used below, private R\&D
investment rose from 2.51\% of GDP in 2019Q4 to 2.86\% in 2025Q4, a nominal
growth rate of 8.6\% a year. At the same time, policy uncertainty has
risen sharply: the EPU index averaged 384 in 2025, almost four times its
1986--2019 mean of 100. How credit conditions and policy uncertainty affect
R\&D investment is therefore a question of current interest. The present boom is led by large, cash-rich firms, for which the
credit channel may be weaker than in the historical average estimated here; we
return to this point in Section~\ref{sec:concl}.

We estimate a Bayesian structural vector autoregression (SVAR) for the U.S.
economy over 1986--2019. Credit supply, uncertainty and aggregate supply
shocks are identified by sign restrictions on output, prices and interest
rates and by restrictions on the relative size of the responses of the excess
bond premium (EBP) of Gilchrist and Zakraj\v{s}ek \cite{gilchrist2012} and the
economic policy uncertainty (EPU) index of Baker, Bloom and Davis
\cite{baker2016}. Following the agnostic approach of Uhlig \cite{uhlig2005}, no
restriction is placed on R\&D investment, so the direction of its response is
estimated rather than assumed. We complement the VAR with local projections
\cite{jorda2005} to examine whether the effects persist, and we compute the
implied path of the R\&D capital stock using the perpetual inventory method of
the BEA R\&D Satellite Account \cite{sliker2007}.

Credit supply shocks lower R\&D investment and account for about 16\% of its
forecast error variance at a five-year horizon, while the effect of
uncertainty shocks is small and not robust. The paper adds to the literature
on R\&D cyclicality, which mostly relates R\&D to output, by separating credit
supply from uncertainty shocks in aggregate data, and it complements
firm-level studies of financing constraints \cite{brown2009,aghion2012}.

% ===========================================================================
\section{Related Literature}
% ===========================================================================

Barlevy \cite{barlevy2007} shows that procyclical R\&D can be privately
optimal when the rents from innovation are dissipated quickly. Ouyang
\cite{ouyang2011} finds that industry R\&D responds procyclically to demand
shocks and countercyclically to supply shocks, and Fabrizio and Tsolmon
\cite{fabrizio2014} find stronger procyclicality where innovations are easily
imitated. At the firm level, Aghion et al.\ \cite{aghion2012} show with French
data that R\&D is countercyclical for unconstrained firms and procyclical for
credit-constrained ones, and Brown, Fazzari and Petersen \cite{brown2009} link
the U.S. R\&D boom of the 1990s to the supply of internal and external equity
finance.

The link between R\&D and long-run growth is central to Schumpeterian and
endogenous growth theory \cite{aghionhowitt1992}. Recent panel evidence for
emerging economies is consistent with it: for Brazil, Russia, China and South
Africa over 1996--2015, Huseynli and Durmu\c{s}kaya \cite{huseynli2026} estimate
a random-effects model with Driscoll--Kraay standard errors and find a positive
and significant association between per capita R\&D expenditure and GDP (an
elasticity of about 1.4), but no significant association between renewable
energy output and GDP. They interpret this as support for Schumpeterian and
endogenous growth models. Their estimates concern long-run levels; the present
paper asks how R\&D responds to cyclical credit and uncertainty shocks, which
determines whether its growth contribution is interrupted in downturns.

Real-options theory predicts that uncertainty delays irreversible investment
\cite{bernanke1983}, and Bloom \cite{bloom2009} and Gulen and Ion
\cite{gulen2016} find that uncertainty and policy uncertainty depress
aggregate and corporate investment. For R\&D the prediction is less clear.
High adjustment costs make R\&D smoother than physical investment
\cite{bloom2007rd}, and firm-level studies find that uncertainty can raise
R\&D \cite{stein2013,atanassov2015}. Caldara et al.\ \cite{caldara2016} show
that financial and uncertainty shocks are difficult to separate, which is why
we identify them jointly. Methodologically, the paper builds on the
sign-restriction approach of \cite{uhlig2005,rubio2010} and on recent work on
inference with local projections \cite{montiel2021}.

% ===========================================================================
\section{Data}\label{sec:data}
% ===========================================================================

\subsection{Sources and sample}
Table~\ref{tab:data} lists the series. All are taken from FRED, except the EBP,
which is published by the Federal Reserve Board \cite{favara2016}. Daily and
monthly series are converted to quarterly frequency by averaging. The
estimation sample is 1986Q1--2019Q4 ($T=136$), with four quarters of 1985 used
as initial observations; the EPU index is not available before 1985. The
baseline sample ends before the pandemic, whose observations are extreme
outliers for a linear VAR \cite{lenza2022}; a robustness check extends the
sample to 2026Q2 with impulse dummies for 2020Q1--2021Q4.

\begin{table}[t]
\caption{Data}
\label{tab:data}
\centering\footnotesize
\begin{tabular}{>{\raggedright}p{2.7cm}>{\raggedright}p{3.0cm}p{1.6cm}}
\hline
Variable & Source (FRED code) & Unit \\
\hline
R\&D investment, nominal & BEA (Y006RC1Q027SBEA) & $100\ln$, real \\
GDP, nominal & BEA (GDP) & $100\ln$, real \\
GDP deflator & BEA (GDPDEF) & $100\ln$ \\
1-year Treasury yield & FRB (DGS1), daily & percent \\
Excess bond premium & FRB website \cite{favara2016}, monthly & p.p. \\
EPU index & Baker et al.\ (USEPUINDXD), daily & $100\ln$ \\
Baa--10y spread & Moody's (BAA10Y), daily & p.p. \\
VIX & CBOE (VIXCLS), daily & $100\ln$ \\
\hline
\multicolumn{3}{l}{\scriptsize Nominal R\&D and GDP are deflated with the GDP deflator.}
\end{tabular}
\end{table}

\subsection{R\&D investment}
The R\&D variable is private fixed investment in research and development
from the National Income and Product Accounts. Because a quarterly price index
for R\&D investment was not available to us, nominal R\&D is deflated with the
GDP deflator. The identifying restrictions imply that the GDP deflator falls
after credit and uncertainty shocks, so this choice, if anything, biases the
estimated decline in real R\&D towards zero; with log nominal R\&D the
response to a credit supply shock after eight quarters is $-1.15\%$, compared
with $-1.01\%$ for the deflated series. Quarterly NIPA R\&D is partly
interpolated from annual survey data, but its quarterly growth rate is not
unusually smooth (first-order autocorrelation 0.27 over 1986--2019, against
0.37 for real GDP).

\subsection{R\&D capital stock}
Following the BEA/NSF R\&D Satellite Account \cite{sliker2007}, the stock of
R\&D capital is constructed by geometric perpetual inventory with the
mid-year convention,
\begin{equation}
K_t=(1-\delta)K_{t-1}+\left(1-\tfrac{1}{2}\delta\right)I_t ,
\label{eq:pim}
\end{equation}
starting in 1947 with $\delta=0.15$. The initial stock is
$K_{1946}=I_{1947}(1-\frac{1}{2}\delta)/(g+\delta)$, where $g$ is the trend
growth rate of real R\&D investment estimated over the early subsample at which
the $R^2$ of a log-linear trend regression reaches its first peak
\cite{sliker2007}. This gives the subsample 1947--1969 and $g=9.45\%$ a year.
By 1986 the initial stock accounts for less than 0.01\% of $K_t$, and the
average growth of $K_t$ over 1986--2019 lies between 4.2\% and 4.4\% a year for
$\delta$ between 0.10 and 0.25.

\subsection{Credit and uncertainty indicators}
The EBP is the component of corporate bond spreads that is not explained by
expected default risk \cite{gilchrist2012}, and it is therefore a better
indicator of credit supply conditions than the Baa--Treasury spread, which also
moves with expected default. The EPU index measures policy-related
uncertainty from newspaper coverage \cite{baker2016}. The Baa--10-year spread
and the VIX are used as alternative indicators.

\subsection{Stationarity}
Table~\ref{tab:ur} reports augmented Dickey--Fuller (ADF) and KPSS tests. For
real GDP, real R\&D and the one-year yield, the ADF test does not reject a unit
root in levels and the KPSS test rejects stationarity, while first differences
are stationary. The EBP, the EPU index and the VIX are stationary in levels.
The GDP deflator is close to being integrated of order two. Following Sims,
Stock and Watson \cite{sims1990}, the VAR is estimated in levels, which keeps
the coefficient estimates consistent in the presence of unit roots and
possible cointegration without imposing a particular cointegration structure.
Impulse responses at long horizons are, however, not consistently estimated in
this case \cite{phillips1998}, which is one reason for examining persistence
with local projections as well.

\begin{table}[t]
\caption{Unit-root tests, 1985Q1--2019Q4. ADF: $t$-statistic (5\% critical
value); KPSS: statistic (5\% critical value 0.46 for differences)}
\label{tab:ur}
\centering\footnotesize
\resizebox{\columnwidth}{!}{\inputtable{tab_unitroot_short}}
\end{table}

% ===========================================================================
\section{Methodology}
% ===========================================================================

\subsection{Reduced-form VAR}
The vector of endogenous variables is
$y_t=(\mathit{GDP}_t, P_t, \mathit{RD}_t, r_t, \mathit{EBP}_t, \mathit{EPU}_t)'$,
where $\mathit{GDP}_t$ is $100$ times the log of real GDP, $P_t$ is $100$ times
the log of the GDP deflator, $\mathit{RD}_t$ is $100$ times the log of real R\&D
investment, $r_t$ is the one-year Treasury yield in percent, $\mathit{EBP}_t$
is the excess bond premium in percentage points and $\mathit{EPU}_t$ is $100$
times the log of the EPU index. Impulse responses of the log variables are
therefore in percent. The reduced form is
\begin{equation}
y_t=c+\sum_{l=1}^{p}A_l y_{t-l}+u_t,\qquad u_t\sim N(0,\Sigma).
\label{eq:var}
\end{equation}
We compare lag lengths $p=1,\dots,8$ on a common sample with the Akaike (AIC),
Schwarz (BIC) and Hannan--Quinn (HQ) criteria,
$\ln|\hat\Sigma_p|+\kappa_T\,pn^2/T$ with $\kappa_T=2$, $\ln T$ and
$2\ln\ln T$, respectively, where $n=6$. The AIC selects two lags and the BIC and
HQ criteria select one. The baseline uses $p=4$, a common choice for
quarterly macroeconomic VARs, and models with two and six lags are reported as
robustness checks. With four lags, multivariate LM tests reject the absence of residual
autocorrelation at orders two to four ($p$-values 0.03, 0.01 and 0.04, but 0.19
at order one), and Jarque--Bera tests reject normality of the GDP, EBP and EPU
residuals; the R\&D residuals are close to normal ($p=0.79$).

We use the standard diffuse prior $p(B,\Sigma)\propto|\Sigma|^{-(n+1)/2}$,
where $B$ collects $c$ and $A_1,\dots,A_p$. The posterior is then
normal--inverse-Wishart \cite{uhlig2005,kilian2017}: $\Sigma\mid y \sim
\mathcal{IW}(\hat U'\hat U,\,T-m)$ and
$\mathrm{vec}(B)\mid\Sigma,y\sim N(\mathrm{vec}(\hat B),\Sigma\otimes(X'X)^{-1})$,
where $\hat B$ and $\hat U$ are the OLS coefficients and residuals, $X$ is the
matrix of regressors and $m$ its number of columns. Draws that imply an
explosive VAR are discarded (39\% of the draws, as expected for a VAR in
levels with roots close to one).

\subsection{Identification}
The structural shocks $\varepsilon_t$ are related to the reduced-form
residuals by $u_t=A\varepsilon_t$ with $AA'=\Sigma$. Any such impact matrix
can be written as $A=PQ$, where $P$ is the lower-triangular Cholesky factor
of $\Sigma$ and $Q$ is an orthogonal matrix. Following Rubio-Ram\'irez,
Waggoner and Zha \cite{rubio2010}, for each draw of $(B,\Sigma)$ from the
posterior we draw $Q$ from the uniform (Haar) distribution, which serves as
the prior over rotations, and keep the draw if the implied impulse responses
satisfy the restrictions in Table~\ref{tab:signs} on impact and one quarter
after the shock. This is repeated until 1{,}000 draws are retained.

\begin{table}[t]
\caption{Identifying restrictions ($h=0,1$)}
\label{tab:signs}
\centering\footnotesize
\begin{tabular}{lcccccc}
\hline
Shock & GDP & $P$ & RD & $r$ & EBP & EPU \\
\hline
Credit supply & $-$ & $-$ & $\cdot$ & $-$ & $+$ & $\cdot$ \\
Uncertainty & $-$ & $-$ & $\cdot$ & $\cdot$ & $\cdot$ & $+$ \\
Aggregate supply & $-$ & $+$ & $\cdot$ & $\cdot$ & $\cdot$ & $\cdot$ \\
\hline
\multicolumn{7}{l}{\scriptsize $\cdot$ = unrestricted. Relative-magnitude restrictions (\ref{eq:rel}) also apply.}
\end{tabular}
\end{table}

A credit supply shock raises the EBP and lowers output and prices, as in
Gilchrist and Zakraj\v{s}ek \cite{gilchrist2012}, whose EBP shocks are also
followed by lower interest rates. Requiring the one-year yield to fall
distinguishes this shock from a monetary policy tightening, which raises both
spreads and short-term rates. An uncertainty shock raises the EPU index and
lowers output and prices, consistent with evidence that uncertainty shocks act
like negative demand shocks \cite{leduc2016}. Since both shocks move output and
prices in the same direction, they are separated by restrictions on the
relative size of the impact responses of the EBP and EPU. Let $a_{ij}$ denote
the impact response of variable $i$ to shock $j$ and $\sigma_i$ the standard
deviation of the reduced-form residual of variable $i$. We require
\begin{equation}
\frac{a_{\mathrm{EBP},\mathrm{cr}}}{\sigma_{\mathrm{EBP}}}>
\frac{a_{\mathrm{EPU},\mathrm{cr}}}{\sigma_{\mathrm{EPU}}},
\qquad
\frac{a_{\mathrm{EPU},\mathrm{unc}}}{\sigma_{\mathrm{EPU}}}>
\frac{a_{\mathrm{EBP},\mathrm{unc}}}{\sigma_{\mathrm{EBP}}},
\label{eq:rel}
\end{equation}
so that a credit supply shock moves the EBP by more standard deviations than
the EPU index, and an uncertainty shock does the opposite. An aggregate supply
shock moves output and prices in opposite directions. The remaining three
shocks are not identified. R\&D investment is not restricted for any shock.
Because pointwise posterior medians do not correspond to a single model
\cite{fry2011}, the historical decomposition is computed for the draw whose
impulse responses are closest to the pointwise medians.

\subsection{Local projections and the R\&D capital stock}
Impulse responses of a VAR at long horizons are extrapolations of its lag
structure. To check persistence, we also estimate lag-augmented local
projections \cite{jorda2005,montiel2021}
\begin{equation}
x_{t+h}-x_{t-1}=\alpha_h+\beta_h\hat\varepsilon_t+\sum_{l=1}^{3}\Gamma_{h,l}\,y_{t-l}+e_{t+h},
\label{eq:lp}
\end{equation}
for $h=0,\dots,28$, where $x_t=\mathit{RD}_t$ and $\hat\varepsilon_t$ is the
standardized credit supply or uncertainty shock from the VAR. To test for a
medium-run effect, the dependent variable is replaced by the average of
$x_{t+h}-x_{t-1}$ over $h=20,\dots,28$ (five to seven years). Standard errors
are computed with the equal-weighted cosine estimator and Student-$t$ critical
values \cite{lazarus2018}. Because local projections at long horizons are
imprecise in samples of this size \cite{kilian2011}, we checked the coverage of
these intervals by simulating 1{,}000 samples from the estimated SVAR: nominal
95\% intervals for the average response over $h=20,\dots,28$ cover the true
value in about 85\% of the samples, so the reported standard errors are, if
anything, too small.

The R\&D capital stock is a fixed linear filter of past investment
(\ref{eq:pim}), so its response is not estimated separately but computed from
the response of R\&D investment in each posterior draw, using the quarterly
version of (\ref{eq:pim}).

% ===========================================================================
\section{Results}
% ===========================================================================

\subsection{Impulse responses}
Fig.~\ref{fig:irf} shows the responses of real R\&D investment, and
Table~\ref{tab:resp} summarizes them. A one-standard-deviation credit supply
shock raises the EBP by 16 basis points and lowers real GDP by 0.16\% and the
one-year yield by 13 basis points on impact. R\&D investment declines
gradually and reaches its trough of $-1.01\%$ after eight quarters, with a 68\%
credible set of $[-1.80,-0.15]$ and a posterior probability of a decline of
0.88. The response is still $-0.49\%$ after five years, but the credible set
then includes zero. Credit supply shocks account for 16\% of the forecast error
variance of R\&D at the five-year horizon.

The uncertainty shock raises the EPU index by about 12.5 log points and
lowers GDP by 0.17\% on impact. Its effect on R\&D is smaller: the median response reaches
$-0.40\%$ after eight quarters, but the credible set is wide and the posterior
probability of a decline is only 0.67. The aggregate supply shock lowers R\&D
by about 0.5\% after two years, again with a credible set that includes zero.
The unrestricted cross-responses are small: the EBP barely moves after an
uncertainty shock, and the EPU index barely moves after a credit supply shock.

\begin{figure*}[t]
\centering
\inputfig{fig_irf_rd}{0.92\textwidth}
\caption{Responses of real R\&D investment (percent) to one-standard-deviation
shocks. Posterior median and 68\% and 90\% credible sets. The response of R\&D
is not restricted.}
\label{fig:irf}
\end{figure*}

\begin{table*}[t]
\caption{Response of real R\&D investment: posterior median [68\% credible set],
posterior probability of a decline at $h=8$, and share of the 20-quarter
forecast error variance}
\label{tab:resp}
\centering\footnotesize
\inputtable{tab_rd_responses}
\end{table*}

\subsection{Historical decomposition}
Fig.~\ref{fig:hd} decomposes the deviation of R\&D investment from its
baseline projection. The identified credit supply shocks explain only a small
part of the large shortfalls in R\&D. Over 2008Q3--2010Q4, R\&D was on average
18.5\% below its baseline; aggregate supply shocks account for about 11
percentage points of this and the unidentified shocks for about 6, while the
contribution of credit supply shocks is close to zero. Over 2001--2003 the
shortfall of 8.6\% is attributed mostly to unidentified shocks. Thus credit
supply shocks lower R\&D on average, but in this decomposition they do not
account for the main R\&D downturns of the sample. This result rests on a
single draw from the identified set and should be read with caution; the identified set also
depends on the uniform prior over rotations \cite{baumeister2015}.

\begin{figure}[t]
\centering
\inputfig{fig_hd_rd_small}{\columnwidth}
\caption{Historical decomposition of real R\&D investment, percent deviation
from the baseline projection (draw closest to the median responses).}
\label{fig:hd}
\end{figure}

\subsection{Persistence}
The local projections in Fig.~\ref{fig:lp} broadly agree with the VAR over the
first two years: after a credit supply shock, R\&D investment falls to about
$-0.7\%$ after nine to ten quarters. Afterwards, however, the local projection
returns to zero after about three years, while the VAR response
decays more slowly. After an uncertainty shock, the local projection is close
to zero for two years and then falls to about $-1\%$ after four years before
returning to zero; this pattern is estimated imprecisely. Table~\ref{tab:lr}
reports the average responses over five to seven years. The local projection
estimates are $-0.20\%$ (s.e.\ 0.31) for the credit supply shock and $0.10\%$
(s.e.\ 0.22) for the uncertainty shock, neither significantly different from
zero. The VAR gives a larger medium-run response to credit supply shocks,
$-0.47\%$, but its credible set includes zero. Because the capital stock
accumulates past investment, the implied R\&D capital stock is $0.46\%$ below
its baseline five to seven years after a credit supply shock, with a 68\%
credible set that just excludes zero. The data therefore do not show a lasting
effect on R\&D investment itself, but the temporary decline is large enough to
leave a persistent gap in the stock of R\&D capital.

\begin{figure}[t]
\centering
\inputfig{fig_lp_rd}{\columnwidth}
\caption{Local projections of real R\&D investment (percent) on the identified
shocks, with 68\% and 90\% confidence intervals.}
\label{fig:lp}
\end{figure}

\begin{table}[t]
\caption{Average response over $h=20,\dots,28$, percent: local projection
(s.e.) and VAR posterior median [68\% credible set]}
\label{tab:lr}
\centering\footnotesize
\inputtable{tab_longrun_short}
\end{table}

\subsection{Robustness}
Table~\ref{tab:rob} reports the response of R\&D after eight quarters in
alternative specifications. The decline of R\&D after a credit supply shock is
robust: its posterior probability lies between 0.74 and 1.00 in all
specifications, and the median response lies between $-0.61\%$ and $-1.44\%$.
The weakest result is obtained when the restriction on the one-year yield is
removed, which makes it harder to separate credit supply from monetary policy
shocks. The response to uncertainty shocks is not robust: it becomes positive
when the VIX replaces the EPU index, and it becomes large and negative only
when the restrictions are imposed for four quarters or when the financial
variables are ordered first in a recursive identification.

\begin{table}[t]
\caption{Robustness: response of R\&D at $h=8$, posterior median [68\% credible
set] and probability of a decline. Percentage points of GDP for R\&D/GDP}
\label{tab:rob}
\centering\footnotesize
\resizebox{\columnwidth}{!}{\inputtable{tab_robustness_short}}
\end{table}

% ===========================================================================
\section{Discussion and Conclusion}\label{sec:concl}
% ===========================================================================

The results support the financing channel of R\&D cyclicality. Tighter credit
supply reduces aggregate U.S. R\&D investment, in line with firm-level
evidence on financing constraints \cite{brown2009,aghion2012}. Once credit
conditions are controlled for, we find no robust effect of policy uncertainty
on R\&D. This is consistent with the view that high adjustment costs make
R\&D less sensitive to uncertainty than physical investment
\cite{bloom2007rd,stein2013}, and with the finding of Caldara et al.\
\cite{caldara2016} that part of what appears to be an uncertainty effect
reflects financial conditions. The effects of credit supply shocks on R\&D
investment fade after three to five years, but they leave a gap in the stock
of R\&D capital. Given the long-run association between R\&D and output
documented in the growth literature \cite{huseynli2026}, such gaps are one way
in which financial disturbances could affect productivity, although testing
this directly is beyond the scope of this paper.

Several limitations should be kept in mind. Aggregate data cannot distinguish
constrained from unconstrained firms. This matters for the current AI-related
investment cycle, which is led by large firms with ample internal funds, for
which the credit channel estimated over 1986--2019 may be weaker. R\&D is
deflated with the GDP deflator. The sample contains few large credit events,
and the results are conditional on the identifying restrictions. Combining the
identified aggregate shocks with firm-level R\&D data would allow the credit
channel to be examined separately for constrained and unconstrained firms.

\section*{Data Availability}
All data are publicly available from the sources in Table~\ref{tab:data}. The
MATLAB code is available from the author on request.

\bibliographystyle{IEEEtran}
\bibliography{references}

\end{document}
